% Additive authoring fragment. It does not modify the current Results chapter.
\subsection{Terminal Robustness and Evidence Synthesis}
\label{subsec:terminal-robustness}

The terminal campaign combines the frozen primary evaluation, post-holdout causal experiments, two retrospective walk-forward windows, synthetic-dose and presentation-budget experiments, and matched pretraining-source controls without pooling their evidential status. All target-weight paths are deterministically rescored on common realised returns. The reconstructed 127,022 daily portfolio observations reconcile to the corresponding monthly gross and net returns to numerical precision, and a clean-room rerun reproduces all terminal tables byte for byte. No policy is retrained and the campaign creates no new confirmatory claim.

Table~\ref{tab:terminal-primary-daily-risk} reports the final primary economic comparison together with the daily mark-to-market audit. At the registered 10-basis-point transaction cost, the NN-vine TD3 ensemble earns an annual CRRA certainty equivalent of 27.96\%, above equal weight, shrinkage mean--variance, rolling vine, and DCC--GARCH, but below the static and dynamic NN-vine optimizers. Its annualized monthly volatility of 13.17\% is below both direct vine optimizers. The result is therefore economically competitive rather than dominant. The benchmark-family White reality-check $p$-value is 0.531 under moving-block resampling and 0.526 under stationary resampling, so the data do not establish a superior strategy after accounting for the searched benchmark family.

Daily reconstruction materially changes the risk interpretation. Figure~\ref{fig:terminal-intr-month-risk} shows that the ensemble's maximum drawdown rises from 6.64\% at monthly endpoints to 13.45\% on the daily path. Equal weight and DCC--GARCH retain lower daily drawdown and 95\% daily CVaR, whereas the static and dynamic NN-vine optimizers exhibit substantially higher daily volatility and tail loss. The ensemble consequently occupies an intermediate downside-risk position. The 99\% tail estimates contain only five observations per locked path and are treated as descriptive sensitivity statistics rather than precise tail-risk evidence.

Figure~\ref{fig:terminal-contrast-forest} synthesizes the registered effects while retaining their evidence classes. The frozen TD3--equal-weight certainty-equivalent difference is positive but imprecise, while the effects relative to static and dynamic NN-vine optimization are close to zero and likewise unresolved. In the post-holdout causal experiment, adding the raw vine state to a policy that already observes scenario CVaR lowers annual certainty equivalent by 7.21 percentage points; all nine block-bootstrap specifications place the interval below zero. Across the two retrospective windows, exposing both the vine state and compressed CVaR observation lowers certainty equivalent by 4.99 points relative to masking both signals. This adverse-direction evidence does not create a new confirmatory winner, but it rejects the interpretation that richer policy-visible dependence information necessarily improves allocation.

The combined evidence instead favors a narrower mechanism. The vine remains active in the synthetic joint-return distribution and scenario-dependent reward even when its fitted features are hidden from the actor. Dependence modelling is therefore more defensible as an implicit training and risk-shaping device than as a high-dimensional policy state. The masked-state policy has the highest descriptive certainty equivalent in both focused windows, but these windows are retrospective and cannot be used for post-holdout model selection.

The matched pretraining controls in Table~\ref{tab:terminal-pretraining-tradeoff} and Figure~\ref{fig:terminal-pretraining-tradeoff} further qualify the synthetic-data contribution. Historical-prefix training has the strongest ensemble point estimate and daily risk profile, whereas concentrated vine-synthetic training is 1.28 certainty-equivalent points lower with a confidence interval spanning economically large effects in both directions. Concentrated synthetic pretraining is 6.54 points above moving-block-bootstrap pretraining, is positive under every leave-one-period-out perturbation, and is better in eight of ten paired seeds; nevertheless, its confidence interval crosses zero and its family-level reality-check $p$-value is approximately 0.24. The defensible conclusion is not that synthetic data raise expected return. Rather, simulation-based pretraining reduces optimization-seed dispersion: the across-seed certainty-equivalent standard deviation falls from 15.65 points under historical-only training to 7.61 under concentrated vine-synthetic training and 4.47 under moving-block bootstrap.

Implementation robustness is favorable to the RL ensemble. Appendix Figure~\ref{fig:terminal-friction} rescoring the same frozen weights shows that TD3 overtakes the dynamic NN-vine optimizer at approximately 12.1 basis points per one-way trade and the static vine at 23.2 basis points. The synthetic masked policy overtakes historical-prefix training at approximately 43.1 basis points because its turnover is lower. These crossings are economic sensitivities, not additional performance trials.

Taken together, the terminal evidence supports a competitive recurrent portfolio policy with useful implementation properties, but not universal benchmark superiority. Its most informative contribution is mechanistic: direct dependence-state injection is unstable, scenario-based dependence can remain useful without being shown to the actor, and synthetic pretraining trades average realized performance against optimization stability. The terminal protocol therefore stops all further same-holdout training and reserves genuine confirmation for independently frozen external or newly arriving data.
